% prop-inner: limit-shapes.tex:903-919 (label prop:inner)
\paragraph{Paper excerpt.}
\begin{verbatim}
\begin{proposition}[Inner radius]\label{prop:inner}
Let $\gauge$ be the Euclidean norm, and let $0<\drift<\nf1d$. For every $p>0$ there exists $C(d,\drift,p)<\infty$ such that, for every integer~$n\geq2$, with probability at least $1-Cn^{-p}$, the following hold simultaneously, where $q_n$ is defined in~\eqref{eq:qn}.
\begin{enumerate}[label=\textup{(\roman*)}]
\item \underline{\emph{Inner radius}}:
\begin{equation}\label{eq:inradius}
|\Rin(n)-r_n|\leq Cr_nq_n\, .
\end{equation}
\item \underline{\emph{Volume}}:
\begin{equation*}
\bigl|r_n^{-1}D_n\mathbin{\triangle}B(0,1)\bigr|\leq Cq_n\, .
\end{equation*}
\item \underline{\emph{Local times}}:
\begin{equation*}
\sup_{y\in\R^d}\bigl|\widetilde\ell_n(y)-2d\drift(r_n-|y|)_+\bigr|\leq Cr_nq_n^{\nf1d}\, .
\end{equation*}
\end{enumerate}
\end{proposition}
\end{verbatim}
\paragraph{Lean statement.}
\begin{verbatim}
theorem CERW.Frozen.inner_radius {d : ℕ} (hd : 2 ≤ d) :
    let ωd : ℝ := (volume (Metric.ball (0 : EuclideanSpace ℝ (Fin d)) 1)).toReal
    ∀ ε : ℝ, 0 < ε → ε < 1 / (d : ℝ) →
    let r : ℕ → ℝ := fun n => ((d + 1) * n / (2 * d * ε * ωd)) ^ ((1 : ℝ) / (d + 1))
    let q : ℕ → ℝ := fun n =>
      if d = 2 then Real.sqrt (Real.log n / r n) else Real.log n / r n
    ∀ p : ℝ, 0 < p → ∃ C : ℝ, 0 < C ∧
      ∀ {Ω : Type u} [MeasurableSpace Ω] (μ : Measure Ω) [IsProbabilityMeasure μ]
        (X : ℕ → Ω → Site d), CERW.IsCERW μ ε X → ∀ n : ℕ, 2 ≤ n →
        μ {ω | ¬ (|CERW.innerRadius (X · ω) n - r n| ≤ C * r n * q n ∧
                  volume (((r n)⁻¹ • CERW.cellSet (X · ω) n) ∆
                      Metric.ball (0 : EuclideanSpace ℝ (Fin d)) 1)
                    ≤ ENNReal.ofReal (C * q n) ∧
                  ∀ y : EuclideanSpace ℝ (Fin d),
                    |CERW.cellLocalTime (X · ω) n y - 2 * d * ε * max (r n - ‖y‖) 0|
                      ≤ C * r n * q n ^ ((1 : ℝ) / d))}
          ≤ ENNReal.ofReal (C * (n : ℝ) ^ (-p))
--
\end{verbatim}
\paragraph{Transcription.}
\emph{Conventions.} \verb!ε! is the paper's $\drift=\kappa$; \verb!hd : 2 ≤ d!, \verb!0 < ε!, \verb!ε < 1 / d! are the standing assumptions $d\ge2$, $0<\drift<1/d$. $\omega_d$ is \verb!(volume (Metric.ball 0 1)).toReal! (as \verb!CERW.unitBallVolume d!). $r_n$ is \verb!((d+1) * n / (2 * d * ε * ωd)) ^ (1/(d+1))! (eq:radius), every cast at a leaf (checked with \verb!pp.coercions.types!). The walk is \verb!CERW.IsCERW μ ε X! on an arbitrary probability space (universe $u$); \verb!X · ω! is the path $n\mapsto X_n(\omega)$. ``With probability at least $1-Cn^{-p}$'' is an upper bound on the outer measure of the exceptional set. \emph{Quantifiers.} $\forall p>0\ \exists C>0\ \forall(\Omega,\mu,X)$ with \verb!IsCERW μ ε X!, $\forall n\ge2$. So $C=C(d,\varepsilon,p)$ precedes the probability space; the paper uses one $C$ for the three bounds and for the probability, and so does the Lean
